% Figure polypeps_encoded_frame.
% Owner: none (note-only figure).
% Used by: gapped-2d-notes PP Sec. V C (Definition 4, encoded frame), Fig. 5(b).
% Source: arXiv OpenAI2026PolynomialPEPS, sections/05-frames.tex:28-82 (inner and
% outer squares, (6.1)-(6.3), Definition 6.1) and sections/06-geometry.tex:109-125
% (true vertices carry holes; the tag holder is one incident owner).
% Ink-to-index: dots are the raw registers of one sheet, one per site. The
% tints are the raw ownership assignment: P_1 (orange, cols 1-5), P_2 (cyan,
% rows 1-5 x cols 6-10), P_3 (grey, rows 6-10 x cols 6-10). The three labels
% meet at the true vertex c (red centre dot, label c above it), the lattice
% point between rows 5-6 and cols 5-6. The hole there is drawn in the
% \holeglyph code (shared-spec 4.1, H5): inner square Q^-(c,h) a solid red
% untinted square of strings at sup-distance h = 1 from c (sites rows 5-6 x
% cols 5-6), outer square Q^+(c,h) a dashed red square of strings at
% sup-distance 2h = 2 (sites rows 4-7 x cols 4-7), both through half-integer
% points, so their sides are in the paper's ratio 1:2; every selected sample
% X_j lies between them.
% Labels: Q^+ is a direct label just east of the outer square's NE corner, in
% the empty middle of four P_2 sites; Q^- ends a red leader that leaves the
% inner square along the white P_2|P_3 gap and ends east of the lattice (the
% gap is too narrow to hold the label without touching the P_2 and P_3
% contours).  The box J (placed between
% sites, inside P_1, label "J (tag)") is the hole's tag register, held by the
% incident owner P_1; its thin red string runs down the white site gap of
% column 1-2 and along the P_1 row gap to the inner square, and only says
% which hole it records.  Along the P_1|P_3 edge only two labels meet, so there
% is no hole: the grey open ring at the bottom end of that edge, labelled
% "no hole".  Raw registers stay at every site, including inside the hole
% (the encoder writes product zeros there).
% Boundary signature: empty (geometry, no tensor contraction).
% Not to scale: drawn with h = 1 lattice spacing; in the paper
% h = h_n = epsilon_0 min(n, t), t = D log L, and outer squares of distinct
% holes are disjoint.
% Related: peps_dyadic_rim_pattern shows how such a true vertex arises after
% point treatment (rim, treated square); this figure has no treatment and no
% rim. peps_frame_hole_encoder_a draws the nested squares alone.
\begin{center}
    \tenkzkernel
    \begin{tenkz}[rows={ket,ket,ket,ket,ket,ket,ket,ket,ket,ket}, cols=10, bonds=none]
        \tndeclare{species}{hole}{hue=source:red}
        \tndeclare{species}{centre}{hue=source:red}
        \tndeclare{species}{bandA}{hue=source:orange}
        \tndeclare{species}{bandB}{hue=source:cyan}
        \tndeclare{species}{bandC}{hue=source:gray}
        \tnmark[form=enclosure, species=bandA, tint, label pos=w]{(1,1) .. (10,5)}{$P_1$}
        \tnmark[form=enclosure, species=bandB, tint, label pos=e]{(1,6) .. (5,10)}{$P_2$}
        \tnmark[form=enclosure, species=bandC, tint, label pos=e]{(6,6) .. (10,10)}{$P_3$}
        \tnwire[kind=string, species=hole, closed, route=orth, name=frInner,
            via={0.5 n of 0.5 w of (5,5), 0.5 n of 0.5 e of (5,6),
                    0.5 s of 0.5 e of (6,6), 0.5 s of 0.5 w of (6,5)}]
        \tn[at=0.5 s of 0.5 e of (5,5), name=frVertex, skin=dot, species=centre, size=l]{}
        \tnmark[form=label, label pos=n]{frVertex}{$c$}
        \tnwire[kind=string, species=hole, stroke=dashed, closed, route=orth, name=frOuter,
            via={0.5 n of 0.5 w of (4,4), 0.5 n of 0.5 e of (4,7),
                    0.5 s of 0.5 e of (7,7), 0.5 s of 0.5 w of (7,4)}]
        \tnmark[form=label, label pos=e]{0.5 n of 0.5 e of (4,7)}{$Q^+$}
        \tnwire[kind=string, species=hole, name=frInnerLead,
            cross={over at crossing of frInnerLead and frOuter}]
        {0.56 e of 0.5 s of (5,6)}{0.6 e of 0.5 s of (5,10)}
        \tnmark[form=label, label pos=e]{0.6 e of 0.5 s of (5,10)}{$Q^-$}
        \tn[at=0.5 s of 0.5 e of (2,1), name=frTag, skin=box]{}
        \tnmark[form=label, label pos=e]{frTag}{$\mathcal J\ (\text{tag})$}
        \tnwire[kind=string, species=hole, route=orth, name=frTagLink,
            cross={over at crossing of frTagLink and frOuter},
            via={0.5 s of 0.5 e of (5,1)}]
        {frTag}{0.56 w of 0.5 s of (5,5)}
        \tn[at=0.5 e of (10,5), name=frNoHole, skin=ring, species=complement]{}
        \tnmark[form=label, label pos=s]{frNoHole}{$\text{no hole}$}
    \end{tenkz}
\end{center}
